Ein «perfekt» isolierter Kasten wird durch eine Wand in zwei gleich grosse Hälften unterteilt. In der einen Hälfte befindet sich ein Gas, die andere Hälfte ist evakuiert. Wird die Wand entfernt, verteilt sich das Gas ohne Arbeit zu verrichten im ganzen Kasten.
\begin{abcliste}
\abc Wie ändert sich die innere Energie des Gases?
\abc Bei realen Gasen beobachtet man bei diesem Versuch eine leichte Temperaturabnahme. Was bedeutet dies für die Zusammensetzung der inneren Energie?
\abc Was geschieht, wenn das Gas durch eine enge Düse in die evakuierte Kammer überströmt und dabei Reibungsarbeit verrichten muss?
\end{abcliste}
